%% 自動生成（tools/mkchecklist.py）。各章の「ふりかえり」の項目を集めたもの
\begingroup\footnotesize
\def\bkcb{\tikz[baseline=-0.5ex]\draw[bkTheme, line width=0.7pt, rounded corners=0.5pt] (0,-0.3em) rectangle (0.8em,0.5em);}
\begin{longtblr}{colspec={Q[c,m,wd=9mm] X[l,m] Q[c,m,wd=10mm] Q[c,m,wd=9mm] Q[c,m,wd=9mm]}, width=\linewidth, rowhead=1, colsep=3pt,
  row{1}={bg=bkVol!14, font=\sffamily\bfseries\scriptsize}, hlines={0.4pt,bkLine}, rowsep=0.7pt}
\R{章}{しょう} & できるように なった こと & ページ & 1\R{回}{かい}め & 2\R{回}{かい}め \\
{\sffamily\bfseries\color{bkTheme}1-1} & 1つずつ ペアに して くらべられる。 & \pageref{bk:1-1:pea} & \bkcb & \bkcb \\
 & かぞえて、 すうじで かける。 & \pageref{bk:1-1:kazoeru} & \bkcb & \bkcb \\
 & おなじ かずは、 おなじ すうじ。 & \pageref{bk:1-1:onaji} & \bkcb & \bkcb \\
{\sffamily\bfseries\color{bkTheme}1-2} & 10ずつ たばねて かぞえられる。 & \pageref{bk:1-2:tabaneru} & \bkcb & \bkcb \\
 & 23の 2と 3の いみが わかる。 & \pageref{bk:1-2:kurai} & \bkcb & \bkcb \\
 & 100、 1000、 10000が わかる。 & \pageref{bk:1-2:hyaku} & \bkcb & \bkcb \\
{\sffamily\bfseries\color{bkTheme}1-3} & なにも ない ことを 0で あらわせる。 & \pageref{bk:1-3:nai} & \bkcb & \bkcb \\
 & 0に あわせて ながさを よめる。 & \pageref{bk:1-3:awaseru} & \bkcb & \bkcb \\
 & なにも ない くらいに 0を かける。 & \pageref{bk:1-3:kurai} & \bkcb & \bkcb \\
{\sffamily\bfseries\color{bkTheme}1-4} & 「あわせる」「ふえる」を しきに かける。 & \pageref{bk:1-4:awaseru} & \bkcb & \bkcb \\
 & ものさしの うえで たし\R{算}{ざん}が できる。 & \pageref{bk:1-4:monosashi} & \bkcb & \bkcb \\
 & 10を こえる たし\R{算}{ざん}が できる。 & \pageref{bk:1-4:kuriagari} & \bkcb & \bkcb \\
{\sffamily\bfseries\color{bkTheme}1-5} & 「のこり」「ちがい」を しきに かける。 & \pageref{bk:1-5:nokori} & \bkcb & \bkcb \\
 & どれだけ かわったかを もとめられる。 & \pageref{bk:1-5:kawatta} & \bkcb & \bkcb \\
 & 10の まとまりを くずして ひける。 & \pageref{bk:1-5:kuzusu} & \bkcb & \bkcb \\
{\sffamily\bfseries\color{bkTheme}1-6} & \R{数直線}{すう|ちょく|せん}で たし\R{算}{ざん}・ひき\R{算}{ざん}が できる。 & \pageref{bk:1-6:susumu} & \bkcb & \bkcb \\
 & 2つの かずの ちがいを みつけられる。 & \pageref{bk:1-6:aida} & \bkcb & \bkcb \\
 & 1\R{目}{め}もりが 10や 100でも よめる。 & \pageref{bk:1-6:ookii} & \bkcb & \bkcb \\
{\sffamily\bfseries\color{bkTheme}1-7} & \R{何十}{なん|じゅう}・\R{何百}{なん|びゃく}の けいさんが できる。 & \pageref{bk:1-7:matomari} & \bkcb & \bkcb \\
 & くらいを そろえて、 ひっ\R{算}{さん}が できる。 & \pageref{bk:1-7:kuraigoto} & \bkcb & \bkcb \\
 & くり\R{上}{あ}がり・くり\R{下}{さ}がりが わかる。 & \pageref{bk:1-7:kuriagari} & \bkcb & \bkcb \\
{\sffamily\bfseries\color{bkTheme}1-8} & 「1\,cm の いくつ\R{分}{ぶん}」で ながさを いえる。 & \pageref{bk:1-8:cm} & \bkcb & \bkcb \\
 & ものさしで cm と mm を よめる。 & \pageref{bk:1-8:yomu} & \bkcb & \bkcb \\
 & ながさの たし\R{算}{ざん}・ひき\R{算}{ざん}が できる。 & \pageref{bk:1-8:keisan} & \bkcb & \bkcb \\
{\sffamily\bfseries\color{bkTheme}1-9} & \R{水}{みず}の かさを、 L・dL・mL で いえる。 & \pageref{bk:1-9:kasa} & \bkcb & \bkcb \\
 & てんびんで、 おもさを くらべられる。 & \pageref{bk:1-9:tenbin} & \bkcb & \bkcb \\
 & おもさを、 g・kg で いえる。 & \pageref{bk:1-9:g} & \bkcb & \bkcb \\
{\sffamily\bfseries\color{bkTheme}1-10} & \mbox{\R{時計}{と|けい}}を よんで、 \R{時}{じ}こくが いえる。 & \pageref{bk:1-10:shikumi} & \bkcb & \bkcb \\
 & \R{時}{じ}こくと \mbox{\R{時間}{じ|かん}}の ちがいが わかる。 & \pageref{bk:1-10:jikan} & \bkcb & \bkcb \\
 & 60で くり\R{上}{あ}がる けいさんが できる。 & \pageref{bk:1-10:rokujuu} & \bkcb & \bkcb \\
{\sffamily\bfseries\color{bkTheme}1-11} & \mbox{1\R{目}{め}もり}が いくつかを \R{見}{み}つけられる。 & \pageref{bk:1-11:hitomemori} & \bkcb & \bkcb \\
 & \mbox{\R{温度計}{おん|ど|けい}}や はかりの \R{目}{め}もりを よめる。 & \pageref{bk:1-11:ondokei} & \bkcb & \bkcb \\
 & \R{目}{め}もりの あいだを 「だいたい」で よめる。 & \pageref{bk:1-11:daitai} & \bkcb & \bkcb \\
{\sffamily\bfseries\color{bkTheme}1-12} & 「1つ\R{分}{ぶん} × いくつ\R{分}{ぶん}」の しきが かける。 & \pageref{bk:1-12:matomari} & \bkcb & \bkcb \\
 & かけ\R{算}{ざん}を おはじきの \R{四角}{し|かく}で あらわせる。 & \pageref{bk:1-12:shikaku} & \bkcb & \bkcb \\
 & 2のだん〜5のだんの \R{九九}{く|く}を つかえる。 & \pageref{bk:1-12:dan} & \bkcb & \bkcb \\
{\sffamily\bfseries\color{bkTheme}1-13} & \R{九九}{く|く}の \R{表}{ひょう}で、 こたえを \R{見}{み}つけられる。 & \pageref{bk:1-13:yomikata} & \bkcb & \bkcb \\
 & \R{表}{ひょう}の きまりを \R{見}{み}つけられる。 & \pageref{bk:1-13:fuekata} & \bkcb & \bkcb \\
 & \R{入}{い}れかえても、 こたえは \R{同}{おな}じと わかる。 & \pageref{bk:1-13:irekae} & \bkcb & \bkcb \\
{\sffamily\bfseries\color{bkTheme}1-14} & \R{分}{わ}ける ことを、 わり\R{算}{ざん}の しきに かける。 & \pageref{bk:1-14:hitoribun} & \bkcb & \bkcb \\
 & わり\R{算}{ざん}の こたえを、 \R{九九}{く|く}で \R{見}{み}つける。 & \pageref{bk:1-14:kuku} & \bkcb & \bkcb \\
 & 0で わる わり\R{算}{ざん}は できないと わかる。 & \pageref{bk:1-14:zero} & \bkcb & \bkcb \\
{\sffamily\bfseries\color{bkTheme}1-15} & あまりの ある わり\R{算}{ざん}が できる。 & \pageref{bk:1-15:amari} & \bkcb & \bkcb \\
 & わり\R{算}{ざん}の こたえを たしかめられる。 & \pageref{bk:1-15:tashikame} & \bkcb & \bkcb \\
 & くり\R{返}{かえ}す ものを わり\R{算}{ざん}で しらべられる。 & \pageref{bk:1-15:shuu} & \bkcb & \bkcb \\
{\sffamily\bfseries\color{bkTheme}1-16} & はんぱを \mbox{\R{小数}{しょう|すう}}で あらわせる。 & \pageref{bk:1-16:shousuu} & \bkcb & \bkcb \\
 & \mbox{\R{数直線}{すう|ちょく|せん}}で \mbox{\R{小数}{しょう|すう}}が よめる。 & \pageref{bk:1-16:suuchokusen} & \bkcb & \bkcb \\
 & \mbox{\R{半分}{はん|ぶん}}を \mbox{\R{分数}{ぶん|すう}}で あらわせる。 & \pageref{bk:1-16:bunsuu} & \bkcb & \bkcb \\
{\sffamily\bfseries\color{bkTheme}1-17} & しらべて かぞえ、 \R{表}{ひょう}に まとめられる。 & \pageref{bk:1-17:hyou} & \bkcb & \bkcb \\
 & ぼうグラフが かけて、 よめる。 & \pageref{bk:1-17:bou} & \bkcb & \bkcb \\
 & グラフの 1\R{目}{め}もりを たしかめられる。 & \pageref{bk:1-17:memori} & \bkcb & \bkcb \\
{\sffamily\bfseries\color{bkTheme}1-18} & \R{表}{ひょう}を、 おれ\R{線}{せん}グラフに かける。 & \pageref{bk:1-18:kaku} & \bkcb & \bkcb \\
 & \R{線}{せん}の むきで、 ふえる・へるが わかる。 & \pageref{bk:1-18:muki} & \bkcb & \bkcb \\
 & かたむきで、 かわる はやさが わかる。 & \pageref{bk:1-18:kyuu} & \bkcb & \bkcb \\
{\sffamily\bfseries\color{bkTheme}1-19} & \R{三角形}{さん|かく|けい}と \R{四角形}{し|かく|けい}を \R{見}{み}わけられる。 & \pageref{bk:1-19:sankaku} & \bkcb & \bkcb \\
 & \R{直角}{ちょっ|かく}や \R{長方形}{ちょう|ほう|けい}・\R{正方形}{せい|ほう|けい}が わかる。 & \pageref{bk:1-19:chokkaku} & \bkcb & \bkcb \\
 & \R{円}{えん}の \R{中心}{ちゅう|しん}と \R{半径}{はん|けい}が わかる。 & \pageref{bk:1-19:en} & \bkcb & \bkcb \\
{\sffamily\bfseries\color{bkTheme}1-20} & \R{角}{かく}は 「まわった \R{大}{おお}きさ」と わかる。 & \pageref{bk:1-20:kaku} & \bkcb & \bkcb \\
 & 1\R{回転}{かい|てん}は 360°、 \R{直角}{ちょっ|かく}は 90°と わかる。 & \pageref{bk:1-20:do} & \bkcb & \bkcb \\
 & \R{分度器}{ぶん|ど|き}で \R{角}{かく}を はかれる。 & \pageref{bk:1-20:bundoki} & \bkcb & \bkcb \\
{\sffamily\bfseries\color{bkTheme}1-21} & \R{広}{ひろ}さを、 マスの \R{数}{かず}で くらべられる。 & \pageref{bk:1-21:masu} & \bkcb & \bkcb \\
 & 1\,cm\textsuperscript{2} を もとに、 \R{面積}{めん|せき}を あらわせる。 & \pageref{bk:1-21:cm} & \bkcb & \bkcb \\
 & \R{長方形}{ちょう|ほう|けい}の \R{面積}{めん|せき}を もとめられる。 & \pageref{bk:1-21:kakezan} & \bkcb & \bkcb \\
\end{longtblr}
\endgroup
